\documentclass[11pt,a4paper]{article}
\usepackage[T1]{fontenc}
\usepackage[utf8]{inputenc}
\usepackage[ngerman]{babel}
\usepackage[margin=2.5cm]{geometry}
\usepackage[hidelinks]{hyperref}
\usepackage{array}
\setlength{\parindent}{0pt}
\setlength{\emergencystretch}{2em}
\renewcommand{\arraystretch}{1.25}
\begin{document}
\shorthandoff{"}
\begin{center}
  {\Large\textbf{Webtechnologien -- Übungsblatt 3}}
\end{center}
\section*{Aufgabe 1: Number-Guessing Game II}
\textbf{Aufgabe:} Das Zahlenratespiel aus Blatt 2 als Webseite mit HTML,
CSS und JavaScript umsetzen: Start-Button, Eingabefeld, Check-Button,
dynamisch erzeugter Exit-Button, maximal drei Versuche. Nur Event Listener
und DOM-Operationen, keine \texttt{prompt}/\texttt{alert}.

\medskip
\textbf{Ergebnis:}
\begin{itemize}
  \item \texttt{index.html}: im \texttt{head} Titel, \texttt{style.css} und
    \texttt{game.js} (mit \texttt{defer}). Im \texttt{body} Überschrift,
    Willkommenstext \texttt{p\#welcome}, \texttt{fieldset\#gameArea} mit
    Legende und Zahlenfeld, darunter \texttt{div\#buttons} mit dem
    deaktivierten \texttt{Check}-Button und dem \texttt{Start Game}-Button.
  \item \texttt{style.css}: \texttt{h1} in Verdana, \texttt{2em},
    \texttt{margin: 15px}; \texttt{p} zentriert mit \texttt{margin: 10px};
    Fieldset mit \texttt{margin: auto} zentriert und zuerst
    \texttt{display: none}; \texttt{input:focus} hellblau
    (\texttt{lightblue}); Check-Button grün und zuerst unsichtbar,
    Start-Button blau, Exit-Button (Klasse \texttt{exitButton}) rot.
    Am Anfang ist nur der Start-Button in der Mitte zu sehen.
  \item \texttt{startGame()}: versteckt Willkommenstext und Start-Button,
    zeigt Fieldset und Check-Button, setzt am Eingabefeld
    \texttt{min=1}, \texttt{max=10}, \texttt{step=1} und
    \texttt{required}, zieht die Zahl mit
    \texttt{Math.floor(Math.random() * 10) + 1}, erzeugt die zwei
    Meldungs-\texttt{div}s mit \texttt{createElement} und setzt sie mit
    \texttt{buttons.before(...)} zwischen Fieldset und Buttons. Der rote
    Exit-Button wird ebenfalls erzeugt und mit \texttt{append} in die
    Button-Leiste gehängt.
  \item \texttt{updateCheckButton()} hängt am \texttt{input}-Event des
    Feldes: \texttt{checkButton.disabled = !guessInput.checkValidity()}.
    Durch die gesetzten Attribute ist nur eine ganze Zahl von 1 bis 10
    gültig.
  \item \texttt{checkGuess()}: richtige Zahl $\rightarrow$
    \glqq Congratulations!\grqq{} mit Anzahl der Versuche und
    \texttt{endGame()}. Falsche Zahl $\rightarrow$ \glqq Too high!\grqq{}
    bzw. \glqq Too low!\grqq, der Tipp kommt mit \texttt{push} in das
    Array \texttt{guesses} und wird mit \texttt{join(", ")} angezeigt.
    Nach drei Versuchen: \glqq Game over!\grqq{} mit der richtigen Zahl,
    \texttt{endGame()} sperrt Check-Button und Eingabefeld.
  \item \texttt{exitGame()}: Fieldset und Check-Button verschwinden,
    Willkommenstext und Start-Button erscheinen wieder,
    \texttt{guesses} wird geleert, die beiden Meldungs-\texttt{div}s und
    der Exit-Button werden mit \texttt{remove()} entfernt. Beim nächsten
    Start werden sie neu erzeugt.
\end{itemize}

\textbf{Hinweis:} Laut Angabe ist der Check-Button nach dem Start aktiv,
laut Teil 2 wird er erst bei einer gültigen Zahl aktiv. Beides ist
umgesetzt: nach dem Start ist er aktiv, ein Klick ohne gültige Zahl zeigt
nur \glqq Please enter a number between 1 and 10.\grqq{} und verbraucht
keinen Versuch. Ab der ersten Eingabe steuert \texttt{checkValidity()} den
Button.

\medskip
\textbf{Test:} In Chrome, \texttt{index.html} direkt geöffnet,
Geheimzahl jeweils 5:
\begin{center}
\begin{tabular}{|>{\raggedright\arraybackslash}p{3.4cm}|>{\raggedright\arraybackslash}p{10.3cm}|}
\hline
\textbf{Eingaben} & \textbf{Anzeige} \\ \hline
\texttt{11} & Check-Button bleibt deaktiviert \\ \hline
\texttt{3}, \texttt{9}, \texttt{5} & \glqq Too low!\grqq{} / \glqq Previous guesses: 3\grqq,
  \glqq Too high!\grqq{} / \glqq Previous guesses: 3, 9\grqq,
  \glqq Congratulations! 5 is correct. You needed 3 attempt(s).\grqq,
  Check-Button deaktiviert \\ \hline
Exit Game & Startzustand, Meldungen und Exit-Button entfernt \\ \hline
\texttt{1}, \texttt{9}, \texttt{2} & \glqq Too low! Game over! The correct number was 5.\grqq,
  \glqq Previous guesses: 1, 9, 2\grqq, Check-Button deaktiviert \\ \hline
\end{tabular}
\end{center}
Zusätzlich geprüft: Verdana mit 32\,px (= 2em) und 15\,px Rand, Farben
blau/grün/rot, hellblaues Feld beim Fokus, Attribute \texttt{1,10,1},
keine Fehler in der Konsole.

\medskip
{\small \textbf{Hilfsmittel:} Claude Code (Claude Opus 5.5) für Erstellung, Prüfung und Dokumentation.}
\end{document}
